% Differential-drive kinematics: pivot / in-place spin / gradual curve.
% Three panels share a single drawing convention so panels can be compared.
%
% Conventions:
%   - Top view, vehicle nose pointing UP (+x_b is forward)
%   - Each track's linear velocity is drawn as a blue arrow whose length
%     is proportional to |v|, with a reference length of REFVEL cm for |v|=v
%   - ICR is a red ringed dot; the upright acronym "ICR" is set with
%     \mathrm{} so it matches the typographic family of v_L, v_R, b
%   - Track width b is dimensioned in panel (a); the same b applies to (b),(c)
%   - Body-frame axes (x_b, y_b) appear once, in panel (b)
\begin{tikzpicture}[
    scale=1.18, every node/.style={transform shape},
    font=\small,
    >={Stealth[length=2.0mm, width=1.5mm]},
    track/.style={draw=black!75, thick, fill=black!22, rounded corners=1pt},
    body/.style={draw=black!85, thick, fill=cppgreen!10, rounded corners=2.5pt},
    nose/.style={draw=black!85, thick, fill=cppgreen!25}, % small heading triangle
    vel/.style={->, line width=1.0pt, codeblue},
    velzero/.style={draw=codeblue, fill=codeblue!30, circle, inner sep=1.3pt},
    icr/.style={circle, draw=warnred, line width=0.85pt, fill=warnred!30, minimum size=0.32cm, inner sep=0pt},
    icrlabel/.style={font=\scriptsize, warnred, inner sep=1pt},
    radlabel/.style={font=\scriptsize, gray!70!black, fill=white, inner sep=1.2pt},
    velabel/.style={font=\scriptsize, codeblue, inner sep=1pt},
    arc/.style={->, semithick, densely dashed, warnred!75},
    title/.style={font=\small\bfseries},
    dim/.style={<->, gray!65, semithick, >={Stealth[length=1.4mm]}},
    extline/.style={gray!50, thin},
    dimtxt/.style={font=\scriptsize, gray!70!black},
    radline/.style={gray!55, thin, densely dashed},
]

% =========================================================================
% Panel (a): PIVOT TURN     v_L = 0,  v_R = +v
% =========================================================================
\begin{scope}[shift={(0,0)}]
    \node[title] at (1.0, 3.05) {(a) Pivot turn};

    % body + tracks
    \draw[body] (0.2,-1.0) rectangle (1.8,1.0);
    % heading nose (triangle on top edge to indicate "front")
    \fill[nose] (0.6,1.0) -- (1.0,1.30) -- (1.4,1.0) -- cycle;
    \draw[track] (-0.1,-1.1) rectangle (0.3,1.1);
    \draw[track] (1.7,-1.1) rectangle (2.1,1.1);

    % --- velocity arrows (proportional: reference length 1.0 cm = |v|) ---
    % left track stationary: blue ringed dot (placed BELOW the ICR on the same track)
    \node[velzero] (lz) at (0.1, -0.55) {};
    \node[velabel, anchor=north] at (0.1, -0.75) {$v_L\!=\!0$};
    % right track: +v forward, length = 1.0 — label CLEAR of the arrow tip
    \draw[vel] (1.9, -0.50) -- (1.9, 0.50);
    \node[velabel, anchor=south west, inner sep=1pt] at (2.05, 0.50) {$v_R\!=\!+v$};

    % --- ICR coincident with stopped left track centerpoint ---
    \node[icr] (icr1) at (0.1, 0.30) {};
    \draw[radline] (icr1.east) -- (1.85, 0.30);  % visualizes that ICR is on left-track axis
    \node[icrlabel, anchor=south] at (0.10, 0.48) {$\mathrm{ICR}$};

    % --- rotation arc (CCW about ICR, sweeping from front to right side) ---
    \draw[arc] (1.85, 1.10) arc[start angle=0, end angle=-50, radius=1.85cm];

    % --- track width b dimension ---
    \draw[extline] (-0.1,-1.45) -- (-0.1,-1.10);
    \draw[extline] ( 2.1,-1.45) -- ( 2.1,-1.10);
    \draw[dim]    (-0.1,-1.32) -- ( 2.1,-1.32);
    \node[dimtxt] at (1.0,-1.45) {track width $b$};
\end{scope}

% =========================================================================
% Panel (b): IN-PLACE ROTATION    v_L = -v, v_R = +v
% =========================================================================
\begin{scope}[shift={(4.6,0)}]
    \node[title] at (1.0, 3.05) {(b) In-place rotation};

    \draw[body] (0.2,-1.0) rectangle (1.8,1.0);
    \fill[nose] (0.6,1.0) -- (1.0,1.30) -- (1.4,1.0) -- cycle;
    \draw[track] (-0.1,-1.1) rectangle (0.3,1.1);
    \draw[track] (1.7,-1.1) rectangle (2.1,1.1);

    % --- velocity arrows: equal magnitude, opposite signs (length = 1.0 each) ---
    \draw[vel] (0.1,  0.50) -- (0.1, -0.50);
    \node[velabel, anchor=north east, inner sep=1pt] at (-0.05, -0.50) {$v_L\!=\!-v$};
    \draw[vel] (1.9, -0.50) -- (1.9,  0.50);
    \node[velabel, anchor=south west, inner sep=1pt] at (2.05,  0.50) {$v_R\!=\!+v$};

    % --- ICR at geometric center; label INSIDE the body so it doesn't cross arrows ---
    \node[icr] (icr2) at (1.0, 0.0) {};
    \node[icrlabel, anchor=south] at (1.0, 0.22) {$\mathrm{ICR}$};

    % --- rotation arc (full-ish CCW around ICR) ---
    \draw[arc] (1.0,0) ++(20:1.20) arc[start angle=20, end angle=160, radius=1.20cm];

    % --- body-frame axes (placed once, in panel b) ---
    \begin{scope}[shift={(-1.05,-1.45)}]
        \draw[->, black!80, thin] (0,0) -- (0.50,0) node[font=\scriptsize, anchor=west, inner sep=1pt]{$y_b$};
        \draw[->, black!80, thin] (0,0) -- (0,0.50) node[font=\scriptsize, anchor=south, inner sep=1pt]{$x_b$};
    \end{scope}
\end{scope}

% =========================================================================
% Panel (c): GRADUAL CURVE    v_L < v_R   (both forward)
% =========================================================================
\begin{scope}[shift={(9.2,0)}]
    \node[title] at (1.0, 3.05) {(c) Gradual curve};

    \draw[body] (0.2,-1.0) rectangle (1.8,1.0);
    \fill[nose] (0.6,1.0) -- (1.0,1.30) -- (1.4,1.0) -- cycle;
    \draw[track] (-0.1,-1.1) rectangle (0.3,1.1);
    \draw[track] (1.7,-1.1) rectangle (2.1,1.1);

    % --- velocity arrows: v_L short (0.40), v_R long (1.00) — proportional ---
    \draw[vel] (0.1, -0.20) -- (0.1,  0.20);
    \node[velabel, anchor=south east, inner sep=1pt] at (-0.05, 0.20) {$v_L$};
    \draw[vel] (1.9, -0.50) -- (1.9,  0.50);
    \node[velabel, anchor=south west, inner sep=1pt] at (2.05, 0.50) {$v_R$};

    % --- ICR to the left of the body on the wheel-axle line (y=0) ---
    \node[icr] (icr3) at (-1.70, 0.0) {};
    \node[icrlabel, anchor=south, inner sep=1pt] at (-1.70, 0.18) {$\mathrm{ICR}$};

    % --- single dashed radius line from ICR to far track edge (no overlap) ---
    \draw[radline] (icr3) -- (2.1, 0.0);

    % --- R dimension: from ICR to body centerline, with tick marks and label ---
    \draw[dim] (-1.70, -0.55) -- (1.0, -0.55);
    \draw[extline] (-1.70, 0.0) -- (-1.70, -0.65);
    \draw[extline] (1.0,   0.0) -- (1.0,   -0.65);
    \node[dimtxt, fill=white, inner sep=1.2pt] at (-0.35, -0.55) {turn radius $R$};

    % --- arc segment indicating the swept curve path ---
    \draw[arc] (icr3) ++(10:2.70) arc[start angle=10, end angle=-15, radius=2.70cm];
\end{scope}

\end{tikzpicture}
